\documentclass[11pt,a4paper]{article}
\usepackage[utf8]{inputenc}
\usepackage[T1]{fontenc}
\usepackage{amsmath, amsfonts, amssymb, bm}
\usepackage{geometry}
\usepackage{mathrsfs}

\geometry{left=2.5cm,right=2.5cm,top=2.5cm,bottom=2.5cm}

\title{\textbf{克拉尼板控制与变厚度薄板振动逆向设计理论}}
\author{}
\date{}

\begin{document}

\maketitle

\section*{一、 变厚度自由边界薄板受迫振动方程推导}

\subsection*{1. 基本假设与几何关系}
基于 Kirchhoff-Love 薄板理论，假设中面内无变形，且变形前垂直于中面的直线在变形后仍为直线且垂直于变形后的中面。设薄板厚度为变量 $h = h(x,y)$，材料密度为 $\rho$，杨氏模量为 $E$，泊松比为 $\nu$。

板的挠度（中面的法向位移）为 $w = w(x,y,t)$。根据几何方程，板的曲率与扭率可表示为：
\begin{equation}
    \kappa_x = -\frac{\partial^2 w}{\partial x^2}, \quad \kappa_y = -\frac{\partial^2 w}{\partial y^2}, \quad \kappa_{xy} = -\frac{\partial^2 w}{\partial x \partial y}
\end{equation}

\subsection*{2. 本构方程与内力}
由于板厚 $h(x,y)$ 是空间坐标的函数，板的抗弯刚度 $D$ 也是变参数：
\begin{equation}
    D(x,y) = \frac{E h^3(x,y)}{12(1-\nu^2)}
\end{equation}
对厚度方向积分，可以得到弯矩 $M_x, M_y$ 和扭矩 $M_{xy}$：
\begin{equation}
\begin{aligned} 
    M_x &= -D(x,y) \left( \frac{\partial^2 w}{\partial x^2} + \nu \frac{\partial^2 w}{\partial y^2} \right) \\ 
    M_y &= -D(x,y) \left( \frac{\partial^2 w}{\partial y^2} + \nu \frac{\partial^2 w}{\partial x^2} \right) \\ 
    M_{xy} &= -D(x,y) (1-\nu) \frac{\partial^2 w}{\partial x \partial y} 
\end{aligned}
\end{equation}

\subsection*{3. 变厚度薄板受迫振动控制方程}
取微元体进行受力分析。设横向分布载荷为 $q(x,y,t)$，惯性力为 $-\rho h(x,y) \frac{\partial^2 w}{\partial t^2}$。由微元体的动量平衡和角动量平衡，可得：
\begin{equation}
    \frac{\partial^2 M_x}{\partial x^2} + 2\frac{\partial^2 M_{xy}}{\partial x \partial y} + \frac{\partial^2 M_y}{\partial y^2} + q(x,y,t) = \rho h(x,y) \frac{\partial^2 w}{\partial t^2}
\end{equation}
代入内力表达式并考虑 $D(x,y)$ 的空间变性：
\begin{equation}
\begin{aligned} 
    & \frac{\partial^2}{\partial x^2} \left[ D \left( \frac{\partial^2 w}{\partial x^2} + \nu \frac{\partial^2 w}{\partial y^2} \right) \right] + 2 \frac{\partial^2}{\partial x \partial y} \left[ D (1-\nu) \frac{\partial^2 w}{\partial x \partial y} \right] \\
    + & \frac{\partial^2}{\partial y^2} \left[ D \left( \frac{\partial^2 w}{\partial y^2} + \nu \frac{\partial^2 w}{\partial x^2} \right) \right] + \rho h \frac{\partial^2 w}{\partial t^2} = q(x,y,t) 
\end{aligned}
\end{equation}
展开并合并同类项，得到变厚度薄板振动控制方程的显式形式：
\begin{equation}
    \nabla^2 \left( D \nabla^2 w \right) - (1-\nu) \left( \frac{\partial^2 D}{\partial x^2}\frac{\partial^2 w}{\partial y^2} - 2\frac{\partial^2 D}{\partial x \partial y}\frac{\partial^2 w}{\partial x \partial y} + \frac{\partial^2 D}{\partial y^2}\frac{\partial^2 w}{\partial x^2} \right) + \rho h \frac{\partial^2 w}{\partial t^2} = q(x,y,t)
\end{equation}
其中 $\nabla^2 = \frac{\partial^2}{\partial x^2} + \frac{\partial^2}{\partial y^2}$。若板厚均匀，$D$ 为常数，方程退化为经典形式 $D\nabla^4 w + \rho h \ddot{w} = q$。

\subsection*{4. 连续问题的算子形式：$(\mathcal{K} - \omega^2 \mathcal{M})W = Q$}
设稳态简谐载荷 $Q(x,y)e^{i\omega t}$ 和响应 $W(x,y)e^{i\omega t}$，线性算子方程为：
\begin{equation}
    (\mathcal{K} - \omega^2 \mathcal{M}) W(x,y) = Q(x,y)
\end{equation}
其中，\textbf{质量算子} $\mathcal{M}$ 定义为：$\mathcal{M}[W] = \rho h(x,y) W$；\\
\textbf{刚度算子} $\mathcal{K}$ 定义为：
\begin{equation}
    \mathcal{K}[W] = \frac{\partial^2}{\partial x^2} \left[ D \left( \frac{\partial^2 W}{\partial x^2} + \nu \frac{\partial^2 W}{\partial y^2} \right) \right] + 2 \frac{\partial^2}{\partial x \partial y} \left[ D (1-\nu) \frac{\partial^2 W}{\partial x \partial y} \right] + \frac{\partial^2}{\partial y^2} \left[ D \left( \frac{\partial^2 W}{\partial y^2} + \nu \frac{\partial^2 W}{\partial x^2} \right) \right]
\end{equation}

\section*{二、 能量泛函方法}

\subsection*{1. 连续系统的稳态能量泛函}
对于简谐受迫振动，系统的总能量泛函 $\Pi(W)$ 为：
\begin{equation}
\begin{aligned} 
    \Pi(W) =& \frac{1}{2} \iint_{\Omega} D(x,y) \left[ (\nabla^2 W)^2 - 2(1-\nu)\left( \frac{\partial^2 W}{\partial x^2}\frac{\partial^2 W}{\partial y^2} - \left(\frac{\partial^2 W}{\partial x \partial y}\right)^2 \right) \right] dxdy \\ 
    &- \frac{1}{2} \omega^2 \iint_{\Omega} \rho h(x,y) W^2 dxdy - \iint_{\Omega} F(x,y) W dxdy 
\end{aligned}
\end{equation}
真实物理场满足变分原理 $\delta \Pi = 0$，此时自由边界条件作为系统的自然边界条件。

\subsection*{2. 空间离散与二次型泛函}
利用基函数 $\mathbf{N}$ 离散化：$W(x,y) = \mathbf{N}(x,y) \mathbf{w}$。泛函退化为离散二次型：
\begin{equation}
    \Pi(\mathbf{w}) = \frac{1}{2} \mathbf{w}^T (\mathbf{K} - \omega^2 \mathbf{M}) \mathbf{w} - \mathbf{w}^T \mathbf{f} = \frac{1}{2} \mathbf{w}^T \mathbf{A} \mathbf{w} - \mathbf{w}^T \mathbf{f}
\end{equation}

\subsection*{3. 梯度寻优与 Hessian 矩阵求解}
求解驻点即令梯度为零：$\nabla_{\mathbf{w}} \Pi(\mathbf{w}) = \mathbf{A} \mathbf{w} - \mathbf{f} = \mathbf{0}$。\\
系统的 \textbf{Hessian 矩阵} 为：$\mathbf{H} = \nabla_{\mathbf{w}}^2 \Pi(\mathbf{w}) = \mathbf{A}$。

\section*{三、 伴随方法数学原理 (Adjoint Method)}

\subsection*{1. 问题定义}
控制方程为：$\mathbf{A}(\mathbf{p}) \mathbf{u} = \mathbf{f}$。目标是最小化目标函数：$\min_{\mathbf{p}} \mathcal{L}(\mathbf{u}, \mathbf{p})$。

\subsection*{2. 梯度推导与困境}
全导数为：
\begin{equation}
    \frac{d\mathcal{L}}{d\mathbf{p}} = \frac{\partial \mathcal{L}}{\partial \mathbf{p}} + \frac{\partial \mathcal{L}}{\partial \mathbf{u}} \frac{d\mathbf{u}}{d\mathbf{p}}
\end{equation}
由控制方程对参数 $\mathbf{p}$ 求导得：$\frac{d\mathbf{A}}{d\mathbf{p}} \mathbf{u} + \mathbf{A} \frac{d\mathbf{u}}{d\mathbf{p}} = \mathbf{0}$，从而 $\frac{d\mathbf{u}}{d\mathbf{p}} = -\mathbf{A}^{-1} \left( \frac{d\mathbf{A}}{d\mathbf{p}} \mathbf{u} \right)$。

\subsection*{3. 伴随法推导}
引入 \textbf{伴随变量} $\bm{\lambda}$，满足伴随方程：
\begin{equation}
    \mathbf{A}^T \bm{\lambda} = \left( \frac{\partial \mathcal{L}}{\partial \mathbf{u}} \right)^T
\end{equation}
最终梯度公式简化为：
\begin{equation}
    \frac{d\mathcal{L}}{d\mathbf{p}} = \frac{\partial \mathcal{L}}{\partial \mathbf{p}} - \bm{\lambda}^T \left( \frac{d\mathbf{A}}{d\mathbf{p}} \mathbf{u} \right)
\end{equation}

\subsection*{4. 算法优势}
\begin{itemize}
    \item \textbf{效率高}：计算梯度仅需额外求解一次伴随方程，与参数维度 $M$ 无关。
    \item \textbf{复用性}：可复用正向求解的矩阵分解结果。
    \item \textbf{节省内存}：无需存储自动微分的计算图。
\end{itemize}

\end{document}